%% reading-question-booklet.tex
%% ---------------------------------------------------------------------------
%% A complete English B HL Paper 2 reading comprehension question booklet:
%% three texts, 37 questions, 40 marks — the shape a real HL paper takes.
%%
%%   tectonic -X compile reading-question-booklet.tex
%%
%% Compiling also writes reading-question-booklet.ans, which markscheme.tex
%% turns into the answer key.
%%
%% The questions refer to reading-text-booklet.tex. Every text there is
%% original; line and paragraph references below match it as it compiles.
%% ---------------------------------------------------------------------------
\documentclass[11pt,a4paper]{article}
\usepackage{ibenglishb}

\ibsetsession{QDHS\,--\,2026}
\ibsetmaxmarks{40}

\begin{document}

% ═══════════════════════════════════════ COVER ═════════════════════════════
\ibcover
  {English B -- Higher level -- Paper 2 -- Reading comprehension}
  {8 May 2026}
  {1 h}
  {\begin{ibinstructions}{Question and answer booklet -- Instructions to students}
     \item Write your session number in the boxes above.
     \item Do not open this examination paper until instructed to do so.
     \item Answer all questions. Each question is allocated \textbf{[1 mark]}
           unless otherwise stated.
     \item Answers must be written within the answer boxes provided.
     \item All answers must be based on the appropriate texts in the
           accompanying text booklet.
     \item The maximum mark for this examination paper is \textbf{[40 marks]}.
   \end{ibinstructions}}
  {10 pages}
  {\ibsessioncode}

% ═══════════════════════════════════════ TEXT A ════════════════════════════
\ibtextheading{A}{The library that lends almost anything}

\ibrubricshort
\ibshortq[(how to use) a wallpaper steamer]%
  {What is Farida Osman explaining to a stranger at the start of the text?}
\ibshortq[(up to) a fortnight]%
  {How long may members keep something they have borrowed?}
\ibshortq[people living within a few streets (of the building)]%
  {Who donated almost all of the items in the collection?}
\ibshortq[a calculation]%
  {According to the text, what persuaded the council to fund the expansion?}

\ibrubriccomplete{\ibparas{3}{4}}
\ibcompleteq[eleven power tools kept behind the desk]%
  {Tomas Iversen started the collection with\dots}
\ibcompleteq[repairs proved expensive]%
  {The scheme to lend musical instruments ended partly because\dots}
\ibcompleteq[belong in a lending collection]%
  {The instruments taught Iversen something about which objects\dots}

\ibturnover
\ibpagebreak

\ibrubricmcq
\begin{ibmcq}[B]{The collection at Brookvale is made up mainly of items that
  were\dots}
  \opt{bought by the council.}
  \opt{given by local people.}
  \opt{lent by other libraries.}
  \opt{left behind by borrowers.}
\end{ibmcq}
\begin{ibmcq}[D]{What does the text suggest about the musical instruments?}
  \opt{They were rarely borrowed.}
  \opt{They were returned damaged.}
  \opt{They were too costly to buy.}
  \opt{They were not suited to lending.}
\end{ibmcq}
\begin{ibmcq}[C]{The advice Iversen gives other librarians is best described as\dots}
  \opt{ambitious.}
  \opt{reluctant.}
  \opt{cautious.}
  \opt{uncertain.}
\end{ibmcq}

% ═══════════════════════════════════════ TEXT B ════════════════════════════
\ibtextheading{B}{Five ways to make a new habit stick}

\ibrubricfindword{\iblines{1}{14}}
\ibfindwordq[surge]{a sudden increase}
\ibfindwordq[dormant]{inactive}
\ibfindwordq[fortnight]{two weeks}
\ibfindwordq[prompted]{reminded}

\ibrubricheadings
\begin{ibmatch}
  \mgap[B] \mgap[D] \mgap[F] \mgap[H]
  \mopt{Reward yourself often}
  \mopt{Attach it to something you already do}
  \mopt{Choose a harder goal}
  \mopt{Never miss twice}
  \mopt{Start before you feel ready}
  \mopt{Do it where someone will notice}
  \mopt{Track every detail}
  \mopt{Keep the record simple}
\end{ibmatch}

\ibturnover
\ibpagebreak

\ibrubricshort
\ibshortq[they went about it in a way that could not last]%
  {According to the writer, why do most people who abandon a new habit fail?}
\ibshortq[to establish that this is something you do]%
  {What does the writer say is the purpose of the first fortnight?}
\ibshortq[attendance]%
  {According to the final section, what are you recording when you mark a
   calendar?}

\ibrubricvocab
\begin{ibmatch}
  \mitem[A]{fragile \ibline{21}}
  \mitem[C]{resist \ibline{26}}
  \mitem[E]{collapse \ibline{27}}
  \mopt{easily broken} \mopt{expensive}
  \mopt{avoid}         \mopt{encourage}
  \mopt{failure}       \mopt{growth}
\end{ibmatch}

% ═══════════════════════════════════════ TEXT C ════════════════════════════
\ibtextheading{C}{An extract from \textit{The Long Way Round}}

\begin{ibtrue}[A, B, D, H]{4}
  \opt{Mirela had been told that the road was being repaired.}
  \opt{Mirela waited at the shelter although she expected no bus.}
  \opt{The walk to the town was uphill for most of the way.}
  \opt{Mirela had made the same journey on foot before.}
  \opt{The dog at the Dennehy place did not recognise her.}
  \opt{Mirela's mother wrote to her about the family shop.}
  \opt{Mirela expected to reach the town later than she had planned.}
  \opt{Mirela did not want to explain her journey to other people.}
\end{ibtrue}

\ibturnover
\ibpagebreak

\ibrubrictf
\ibtfq[F: ``she had known it would not'']%
  {Mirela was surprised when the bus failed to arrive.}
\ibtfq[F: ``came to the gate and did not bark'']%
  {The dog at the Dennehy place barked as she went past.}
\ibtfq[T: ``accurate about the surface and silent about everything underneath'']%
  {Mirela thought her mother's letters left a great deal unsaid.}
\ibtfq[F: ``how little she intended to explain'']%
  {Mirela planned to describe her journey to the people of the town.}

\ibrubricrefer
\ibreferq[the bus]{\dots she had known \ibu{it} would not. \ibline{1}}
\ibreferq[the walk (to the town)]{\dots she had done \ibu{it} before\dots \ibline{7}}
\ibreferq[the letters]{\ibu{They} never mentioned her father's brother\dots \ibline{16}}
\ibreferq[the people of the town]{Let \ibu{them} see her come up the hill\dots \ibline{21}}

\ibturnover
\ibpagebreak

\ibrubricmcq
\begin{ibmcq}[C]{Mirela stood at the bus shelter because\dots}
  \opt{she thought the driver might have been wrong.}
  \opt{she wanted to speak to the other passengers.}
  \opt{leaving the house made the day feel normal.}
  \opt{the road had only just been closed.}
\end{ibmcq}
\begin{ibmcq}[B]{``a kind of weather report'' \ibline{15} suggests that the
  letters were\dots}
  \opt{written in a hurry.}
  \opt{true but shallow.}
  \opt{full of complaint.}
  \opt{difficult to follow.}
\end{ibmcq}
\begin{ibmcq}[D]{Why is Mirela not sorry that the lorry drove past?}
  \opt{She had already rested at the crossroads.}
  \opt{She did not recognise the driver.}
  \opt{She knew she would arrive early anyway.}
  \opt{She wanted to arrive on foot.}
\end{ibmcq}
\begin{ibmcq}[A]{In this passage, Mirela is best described as\dots}
  \opt{self-possessed.}
  \opt{resentful.}
  \opt{anxious.}
  \opt{carefree.}
\end{ibmcq}

\ibblankpage

\end{document}
